\chapter{Acotación de sucesiones convergentes}

% ============================================================
\section{Sumas Finitas: Cada Término Acotado por el Total}
% ============================================================

Al trabajar con sumas finitas de cantidades no negativas, emerge una
propiedad fundamental: ningún término individual puede superar la suma
total. Para sucesiones de números reales, esto se traduce en que el
valor absoluto de cualquier término está acotado por la suma de todos
los valores absolutos hasta ese punto.

Este resultado es especialmente útil cuando necesitamos controlar
términos individuales usando información global sobre sumas. La
demostración emplea inducción matemática, lo cual es natural dado que
trabajamos con sumas sobre rangos finitos.

% ────────────────────────────────────────────────────────────
\subsection{Nuevas Herramientas}
% ────────────────────────────────────────────────────────────

\subsubsection{Notación de suma: \texttt{range}}

La notación $\sum_{k \in \mathrm{range}\, N}$ significa sumar con $k$
yendo de $0$ a $N-1$ (lo que da $N$ términos en total). Esta es la
forma estándar de expresar sumas finitas en matemática formal.

\subsubsection{Sucesor del rango de suma: \texttt{sum\_range\_succ}}

Para una función $f : \N \to \R$ y un número natural $N$, se tiene:
\[
  \sum_{n \in \mathrm{range}(N+1)} f(n)
  \;=\;
  \sum_{n \in \mathrm{range}\, N} f(n) + f(N).
\]
Esto permite separar el último término de una suma, lo cual es
esencial en las demostraciones por inducción.

\subsubsection{No negatividad de sumas: \texttt{sum\_nonneg}}

Si una función es no negativa en todas partes, entonces su suma
finita también es no negativa. Más precisamente, si $f(k) \ge 0$
para todo $k$ en el rango, entonces $\sum_k f(k) \ge 0$.

\subsubsection{Prueba por contradicción: \texttt{contradiction}}

Si las hipótesis conducen a una contradicción lógica (por ejemplo, si
una hipótesis afirma $n < 0$ siendo $n : \N$ un número natural),
entonces la táctica \texttt{contradiction} cierra el objetivo
automáticamente.

% ────────────────────────────────────────────────────────────
\subsection{El Enunciado Matemático}
% ────────────────────────────────────────────────────────────

\begin{theorem}[\texttt{TermLeSum}]
  Sea $a : \N \to \R$ una sucesión y $N$ un número natural. Entonces,
  para todo $n < N$:
  \[
    \abs{a(n)} \;\le\; \sum_{k=0}^{N-1} \abs{a(k)}.
  \]
\end{theorem}

% ────────────────────────────────────────────────────────────
\subsection{Enfoque Estratégico}
% ────────────────────────────────────────────────────────────

La idea clave es usar inducción sobre $N$ desde el principio, en
lugar de introducir $n$ primero. Esto puede parecer contraintuitivo
dado que el enunciado comienza con $\forall n < N$, pero la
estructura inductiva hace la demostración mucho más limpia.

Para el \textbf{caso base} $N = 0$, no existe ningún número natural
$n < 0$, por lo que el enunciado es vacuamente verdadero (lo que
conduce a una contradicción al intentar suponer que tal $n$ existe).

Para el \textbf{paso inductivo}, suponiendo que el resultado vale
para $N$, lo probamos para $N + 1$ distinguiendo dos casos: o bien
$n < N$ (donde usamos la hipótesis inductiva), o bien $n = N$ (donde
usamos la no negatividad de las sumas).

% ────────────────────────────────────────────────────────────
\subsection{Solución en Lean}
% ────────────────────────────────────────────────────────────

\begin{minted}{lean}
Statement TermLeSum (a : N → R) (N : N) :
    ∀ n < N, |a n| ≤ ∑ k ∈ range N, |a k| := by
  induction' N with N hN
  · intro n hn
    contradiction
  · intro n hn
    have f1 : ∑ k ∈ range (N + 1), |a k| =
              (∑ k ∈ range N, |a k|) + |a N| := by
      apply sum_range_succ
    rewrite [f1]
    by_cases hn' : n < N
    · specialize hN n hn'
      have f1' : 0 ≤ |a N| := by bound
      linarith [f1', hN]
    · have f2 : n = N := by bound
      rewrite [f2]
      have f3 : ∀ k ∈ range N, 0 ≤ |a k| := by bound
      have f4 : 0 ≤ ∑ k ∈ range N, |a k| := by
        apply sum_nonneg f3
      linarith [f4]
\end{minted}

% ────────────────────────────────────────────────────────────
\subsection{Demostración en Lenguaje Natural}
% ────────────────────────────────────────────────────────────

\begin{proof}
Procedemos por inducción sobre $N$.

\medskip
\noindent\textbf{Caso base.} Si $N = 0$, no existe ningún número
natural $n < 0$, por lo que el enunciado es vacuamente verdadero.

\medskip
\noindent\textbf{Paso inductivo.} Supongamos que el resultado vale
para algún $N \ge 0$, y sea $n < N + 1$. Observamos primero que:
\[
  \sum_{k=0}^{N} \abs{a(k)}
  = \sum_{k=0}^{N-1} \abs{a(k)} + \abs{a(N)}.
\]
Distinguimos dos casos:

\medskip
\noindent\textit{Caso 1: $n < N$.} Por la hipótesis inductiva,
$\abs{a(n)} \le \sum_{k=0}^{N-1} \abs{a(k)}$.
Como $\abs{a(N)} \ge 0$, obtenemos:
\[
  \abs{a(n)}
  \;\le\; \sum_{k=0}^{N-1} \abs{a(k)}
  \;\le\; \sum_{k=0}^{N-1} \abs{a(k)} + \abs{a(N)}
  = \sum_{k=0}^{N} \abs{a(k)}.
\]

\noindent\textit{Caso 2: $n \not< N$ pero $n < N+1$, luego $n = N$.}
Debemos mostrar que $\abs{a(N)} \le \sum_{k=0}^{N} \abs{a(k)}$.
Como $\abs{a(k)} \ge 0$ para todo $k \in \{0,\ldots,N-1\}$, tenemos
$\sum_{k=0}^{N-1} \abs{a(k)} \ge 0$, y por tanto:
\[
  \abs{a(N)}
  \;\le\; \sum_{k=0}^{N-1} \abs{a(k)} + \abs{a(N)}
  = \sum_{k=0}^{N} \abs{a(k)}.
\]
Esto completa la inducción.
\end{proof}

% ────────────────────────────────────────────────────────────
\subsection{Por Qué Importa Este Resultado}
% ────────────────────────────────────────────────────────────

Este resultado ilustra un principio fundamental: los términos
individuales quedan controlados por las sumas agregadas. Esto es
crucial cuando necesitamos acotar términos específicos usando
información global sobre la sucesión. El teorema se aplicará
inmediatamente para demostrar que las sucesiones convergentes son
acotadas.

% ============================================================
\section{Acotamiento de Sucesiones Convergentes}
% ============================================================

Una de las propiedades más importantes de las sucesiones convergentes
es que deben ser acotadas. Intuitivamente, esto tiene sentido: si una
sucesión se aproxima a algún límite finito, no puede alejarse hasta el
infinito. Los términos deben permanecer a una distancia fija del
límite.

Este teorema tiene implicaciones profundas. Nos dice que la
convergencia impone restricciones fuertes sobre una sucesión: las
sucesiones convergentes no pueden exhibir comportamientos salvajes e
ilimitados. Esto se usa extensamente en análisis para descartar ciertos
casos patológicos.

% ────────────────────────────────────────────────────────────
\subsection{Nuevas Herramientas}
% ────────────────────────────────────────────────────────────

\subsubsection{Definición: \texttt{SeqBdd}}

Una sucesión $a : \N \to \R$ es \emph{acotada} (denotado
$\mathrm{SeqBdd}(a)$) si existe algún número real positivo $M$ tal
que $\abs{a(n)} \le M$ para todo $n \in \N$. Formalmente:
\[
  \mathrm{SeqBdd}(a)
  \;\Longleftrightarrow\;
  \exists M > 0,\; \forall n \in \N,\; \abs{a(n)} \le M.
\]

% ────────────────────────────────────────────────────────────
\subsection{El Enunciado Matemático}
% ────────────────────────────────────────────────────────────

\begin{theorem}[\texttt{BddOfConvNonzero}]
  Si $a : \N \to \R$ converge a $L$ con $L \ne 0$, entonces $a$ es
  acotada.
\end{theorem}

% ────────────────────────────────────────────────────────────
\subsection{Enfoque Estratégico}
% ────────────────────────────────────────────────────────────

La estrategia de demostración consta de dos partes:

\begin{enumerate}
  \item \textbf{Eventualmente acotada.} Usamos la convergencia para
    mostrar que, eventualmente (para $n \ge N$), todos los términos
    permanecen a distancia $\abs{L}$ de $L$, y por tanto
    $\abs{a(n)} \le 2\abs{L}$ para $n$ grande.

  \item \textbf{Inicialmente acotada.} Para los finitos términos
    anteriores a $N$, usamos el Teorema \texttt{TermLeSum} para
    acotar cada uno por la suma $\sum_{k=0}^{N-1} \abs{a(k)}$.

  \item \textbf{Cota global.} Combinamos ambas cotas para obtener:
    \[
      M = 2\abs{L} + \sum_{k=0}^{N-1} \abs{a(k)}.
    \]
\end{enumerate}

Esta cota funciona para todos los términos: los términos iniciales
quedan cubiertos por la suma finita, y los términos tardíos quedan
cubiertos por $2\abs{L}$.

% ────────────────────────────────────────────────────────────
\subsection{Solución en Lean}
% ────────────────────────────────────────────────────────────

\begin{minted}{lean}
Statement BddOfConvNonzero (a : N → R) (L : R)
    (ha : SeqLim a L) (hL : L ≠ 0) :
    SeqBdd a := by
  choose N hN using EventuallyBdd_of_SeqConv a L ha hL
  use 2 * |L| + ∑ k ∈ range N, |a k|
  have absL : 0 < |L| := by apply abs_pos_of_nonzero hL
  have f1 : ∀ k ∈ range N, 0 ≤ |a k| := by bound
  have f2 : 0 ≤ ∑ k ∈ range N, |a k| := by apply sum_nonneg f1
  split_ands
  · linarith [f2, absL]
  · intro n
    by_cases hn : N ≤ n
    · specialize hN n hn
      linarith [hN, f2]
    · have hn' : n < N := by bound
      have f3 : |a n| ≤ ∑ k ∈ range N, |a k| := by
        apply TermLeSum a N n hn'
      linarith [f3, absL]
\end{minted}

% ────────────────────────────────────────────────────────────
\subsection{Demostración en Lenguaje Natural}
% ────────────────────────────────────────────────────────────

\begin{proof}
Como $a(n) \to L$ y $L \ne 0$, por el teorema
\texttt{EventuallyBdd\_of\_SeqConv}, existe $N$ tal que
$\abs{a(n)} \le 2\abs{L}$ para todo $n \ge N$.

Definimos:
\[
  M := 2\abs{L} + \sum_{k=0}^{N-1} \abs{a(k)}.
\]

Primero verificamos que $M > 0$. Como $L \ne 0$, tenemos $\abs{L} > 0$,
luego $2\abs{L} > 0$. Además, como $\abs{a(k)} \ge 0$ para cada $k$,
tenemos $\sum_{k=0}^{N-1} \abs{a(k)} \ge 0$. Por tanto $M > 0$.

Ahora mostramos que $\abs{a(n)} \le M$ para todo $n \in \N$.
Consideramos dos casos:

\medskip
\noindent\textit{Caso 1: $n \ge N$.}
Por nuestra elección de $N$:
\[
  \abs{a(n)} \;\le\; 2\abs{L}
  \;\le\; 2\abs{L} + \sum_{k=0}^{N-1} \abs{a(k)} = M.
\]

\noindent\textit{Caso 2: $n < N$.}
Por el Teorema \texttt{TermLeSum}:
\[
  \abs{a(n)} \;\le\; \sum_{k=0}^{N-1} \abs{a(k)}
  \;<\; 2\abs{L} + \sum_{k=0}^{N-1} \abs{a(k)} = M.
\]

Por tanto, $a$ es acotada.
\end{proof}

% ────────────────────────────────────────────────────────────
\subsection{Importancia del Acotamiento}
% ────────────────────────────────────────────────────────────

El acotamiento de las sucesiones convergentes es fundamental en
análisis. Se utiliza para demostrar:

\begin{itemize}
  \item El teorema de Bolzano-Weierstrass (toda sucesión acotada tiene
    una subsucesión convergente).
  \item Diversos resultados de compacidad.
  \item Teoremas de convergencia uniforme.
  \item Propiedades de funciones continuas en conjuntos compactos.
\end{itemize}

Este teorema también contribuye a desarrollar la intuición: converger
hacia un límite finito implica que la sucesión no puede escapar al
infinito. Debe permanecer atrapada en una región acotada de la recta
real.

%%% Local Variables:
%%% mode: LaTeX
%%% TeX-master: "Analisis_en_Lean"
%%% End:
